\section{Exact equivariant quotients and tensor weight error in the
original zeta
receiver}\label{exact-equivariant-quotients-and-tensor-weight-error-in-the-original-zeta-receiver}

Independent mathematical derivation, 24 September 2026. This note uses
the actual Mellin quotient and prime actions defined in IH2--IH3 and
PL6. It proves the quotient maps and their effects rather than treating
``quotientable'' as a synonym for ``pure.'' It does not assume a Deligne
weight estimate for this receiver.

Source notation is retained: \(Z_0\) is absence; \(Z_1\) is primitive
presence \(\tau\), without \(Z_2\) parity. Addition on source \(\tau\)
is not defined here. The vector-space zero and scalar identities below
belong to receiving algebras. Neither is identified with source absence
or with \(\tau\). Every support coefficient is retained by an explicitly
specified identity map in QS8.

\subsection{QS1. The actual local and global
objects}\label{qs1.-the-actual-local-and-global-objects}

Let \(\mathcal A\), \(I\), and \(Q=\mathcal A/I\) be precisely IH2.1: \[
p_{N,j}(k)=\sup_{u>0}(u^N+u^{-N})|(u\partial_u)^jk(u)|<\infty,\qquad
I=\overline{\mathcal E(S_0^{\mathrm{even}})}^{\mathcal A},
\] \[
\mathcal E f(u)=u^{1/2}\sum_{n\ge1}f(nu),\qquad
F_k(s)=\int_0^\infty k(u)u^{s-1/2}\frac{du}{u}.
\tag{QS1.1}
\] The conditions defining \(S_0^{\mathrm{even}}\) are even Schwartz
regularity, \(f(0)=0\), and \(\int_{\mathbb R}f=0\). Every symbol below
is an observation in this receiver of the original, uncompleted zeta
function.

For an actual nontrivial zero \(\rho\) with multiplicity \(m\), retain
\[
A_{\rho,m}=\mathbb C[t_\rho]/(t_\rho^m),\qquad
j_{\rho,m}([k])=\sum_{j=0}^{m-1}\frac{F_k^{(j)}(\rho)}{j!}t_\rho^j.
\tag{QS1.2}
\] For \(a>0\), the clock operator and its full local multiplier are \[
W_a k(u)=a^{1/2}k(u/a),\qquad
w_{a,\rho,m}=a^\rho\sum_{j=0}^{m-1}\frac{(\log a)^j}{j!}t_\rho^j.
\tag{QS1.3}
\] Thus \(j_{\rho,m}W_a=M_{w_{a,\rho,m}}j_{\rho,m}\). The factor
\(a^{1/2}\) in the original operator and every nilpotent coefficient are
retained. IH2--IH3 prove that these maps descend through \(I\), that
\(j_{\rho,m}\) is onto, and that their product is onto for any finite
set of distinct zeros.

No deformation of the arithmetic counting measure is presumed merely
from the existence of a spectral character.

\subsection{QS2. Complete classification on one primary
block}\label{qs2.-complete-classification-on-one-primary-block}

Fix \(p>1\), put \(\lambda=p^\rho\), and write \(T=M_{t_\rho}\). Then \[
M_w=\lambda\exp((\log p)T),\qquad T^m=0,
\] \[
X=\lambda^{-1}M_w-I,\qquad
T=\frac1{\log p}\sum_{j=1}^{m-1}\frac{(-1)^{j+1}}jX^j.
\tag{QS2.1}
\] These are finite polynomial identities: substitute into
\(\log(\exp z)=z\) modulo \(z^m\). Conversely
\(M_w=\lambda\sum_{j=0}^{m-1}(\log p)^jT^j/j!\).

A complex-linear subspace \(K\subset A_{\rho,m}\) is \(M_w\)-invariant
if and only if it is \(T\)-invariant, by these two polynomial
identities. Such a subspace is an ideal, because it is invariant under
every polynomial in \(t_\rho\).

Every ideal is exactly \((t_\rho^r)\) for a unique
\(r\in\{0,\ldots,m\}\). For a nonzero ideal, take its least nonzero
degree \(r\), and select \(t_\rho^r u(t_\rho)\) with \(u(0)\ne0\). The
inverse of \(u\) is a finite geometric polynomial in its nilpotent part,
so the ideal contains \(t_\rho^r\). Minimality of \(r\) says that every
other element is divisible by \(t_\rho^r\). The zero ideal corresponds
to \(r=m\).

Consequently every equivariant quotient is, up to equivariant linear
isomorphism, the exact truncation \[
q_r:A_{\rho,m}\longrightarrow A_{\rho,r}
 =\mathbb C[t_\rho]/(t_\rho^r),\qquad 0\le r\le m.
\tag{QS2.2}
\] Here \(r=0\) denotes the zero vector space; for every \(r\ge1\) the
surviving operator is multiplication by \[
p^\rho\sum_{j=0}^{r-1}\frac{(\log p)^j}{j!}t_\rho^j,\qquad
\det(XI-M_{w,r})=(X-p^\rho)^r.
\tag{QS2.3}
\] A nonzero quotient cannot change \(p^\rho\). It can shorten the
nilpotent block or kill the entire block. These are all possibilities in
this category.

The intertwiner space between blocks at the same \(\rho\), of lengths
\(m,r\), is explicit. Formula (QS2.1) shows that an intertwiner commutes
with \(t_\rho\). It is determined by the image \(v\) of the constant
polynomial and must satisfy \(t_\rho^mv=0\). Thus \[
\operatorname{Hom}_{W_p}(A_{\rho,m},A_{\rho,r})
\cong t_\rho^{\max(r-m,0)}A_{\rho,r}.
\tag{QS2.4}
\] The map sends \(f(t_\rho)\) to \(f(t_\rho)v\); it is surjective
precisely when \(m\ge r\) and \(v(0)\ne0\).

\subsection{QS3. Two counted primes separate all distinct zero
characters}\label{qs3.-two-counted-primes-separate-all-distinct-zero-characters}

For complex \(\rho,\nu\), \[
2^\rho=2^\nu,\quad3^\rho=3^\nu
\quad\Longrightarrow\quad \rho=\nu.
\tag{QS3.1}
\] Indeed \((\rho-\nu)\log2=2\pi ia\) and \((\rho-\nu)\log3=2\pi ib\)
for integers \(a,b\). A nonzero difference forces \(a,b\ne0\) and
\(\log2/\log3=a/b\), contradicting unique factorization after
exponentiation.

For a finite set \(S\) of distinct actual zeros, put \[
B_S=\prod_{\rho\in S}A_{\rho,m_\rho}.
\] There is \(c\in\mathbb C\) for which \(\lambda_\rho=2^\rho+c3^\rho\)
are pairwise distinct. Each equality forbids at most one \(c\), unless
its coefficients of \(c\) agree; in that case (QS3.1) makes the constant
coefficients different, so no \(c\) is forbidden.

Set \(M=M_{w_2}+cM_{w_3}\). On the \(\rho\)-block,
\((M-\lambda_\rho I)^{m_\rho}=0\). The pairwise coprime polynomials
\((X-\lambda_\rho)^{m_\rho}\) have Chinese-remainder projectors
\(P_\rho(X)\): \[
P_\rho\equiv1\bmod(X-\lambda_\rho)^{m_\rho},\qquad
P_\rho\equiv0\bmod(X-\lambda_\nu)^{m_\nu}\quad(\nu\ne\rho).
\] Bezout's identity constructs them. Evaluating at \(M\) gives the
exact projection onto each block.

Any subspace \(K\subset B_S\) invariant under \(W_2,W_3\) is invariant
under these projections. QS2 therefore gives \[
K=\prod_{\rho\in S}(t_\rho^{r_\rho}),\qquad
0\le r_\rho\le m_\rho.
\tag{QS3.2}
\] Conversely every displayed \(K\) is invariant under every \(W_a\).
This is the complete finite spectral quotient classification for all
prime actions.

No nonzero simultaneous \(W_2,W_3\) intertwiner connects blocks at
distinct \(\rho,\nu\). At least one prime has distinct eigenvalues by
(QS3.1); Bezout's identity for its two coprime minimal-polynomial
factors forces that intertwiner to vanish. Thus equivariant isomorphism
cannot transport an off-critical block to a critical block.

This is for the original fixed actions. PL2 treats a simultaneous change
of clock separately: \(\zeta(s)\) becomes \(\zeta(cs)\), and both its
zero and reflection line rescale. That comparison does not make an
off-critical position critical relative to the rescaled line.

\subsection{QS4. Reflection-compatible quotients retain the off-critical
value
form}\label{qs4.-reflection-compatible-quotients-retain-the-off-critical-value-form}

Take the full actual off-critical orbit \[
\mathcal O_\rho=(\rho,\ 1-\overline\rho,\ \overline\rho,\ 1-\rho),
\quad \rho=\beta+i\gamma,\quad0<\beta<1,\quad
\beta\ne\tfrac12,\quad\gamma\ne0.
\] All four multiplicities are \(m\). Reflection is \[
(x^\#)_\omega(t)=\overline{x_{1-\overline\omega}}(-t).
\tag{QS4.1}
\] A quotient in QS3 carries this involution exactly when \[
r_\omega=r_{1-\overline\omega}.
\tag{QS4.2}
\] The original conjugation symmetry further imposes
\(r_\omega=r_{\overline\omega}\). All four lengths then agree. Length
zero kills the whole orbit; any positive length retains all four
characters.

Retain the original-multiplicity form \[
H_m(x,y)=m\sum_{\omega\in\mathcal O_\rho}
\overline{x_{1-\overline\omega}(0)}y_\omega(0).
\tag{QS4.3}
\] Its radical is the product of the four ideals \((t_\omega)\), because
its value matrix is the invertible matrix
\(m\operatorname{diag}(J_2,J_2)\), with \[
J_2=\begin{pmatrix}0&1\\1&0\end{pmatrix}.
\] Every quotient with all \(r_\omega\ge1\) has kernel in that radical,
so the form descends. Its value signature remains \((2,2)\); values
\((1,-1,0,0)\) give \(-2m\). Quotienting higher jets cannot make this
form positive.

The descended form still has coefficient \(m\), the original
multiplicity. Intrinsic trace on a shortened block would have
coefficient \(r_\omega\). Replacing \(m\) by \(r_\omega\) would alter
the original trace datum.

No prescribed finite value vector is asserted to vanish at all other
zeta zeros. IH3 gives surjectivity onto this finite orbit; this
signature statement is about its exact finite receiving form. The global
form retains all other zeros.

\subsection{QS5. The greatest globally invisible quotient is
explicit}\label{qs5.-the-greatest-globally-invisible-quotient-is-explicit}

For the set \(Z\) of distinct actual nontrivial zeros define \[
J_{\mathrm{val}}=\bigcap_{\rho\in Z}\ker(F_\bullet(\rho)),\qquad
J_{\mathrm{jet}}=\bigcap_{\rho\in Z}\ker j_{\rho,m_\rho}.
\tag{QS5.1}
\] These are closed invariant convolution ideals. Closedness follows
from continuity of Mellin derivatives. The ideal property follows from
the character and algebra-map product laws. Full multiplier formulas
prove invariance under \(W_a\); permutation of zeros and jets proves
invariance under the original symmetries.

For \(\pi_K:Q\to Q/K\), all original value observations factor through
\(\pi_K\) exactly when \[
K\subset J_{\mathrm{val}},
\tag{QS5.2}
\] and all full original jets factor exactly when \[
K\subset J_{\mathrm{jet}}.
\tag{QS5.3}
\] Proof: a linear map factors through a quotient exactly when it
vanishes on its kernel. Apply this to each specified observation.

Thus \(Q/J_{\mathrm{val}}\) and \(Q/J_{\mathrm{jet}}\) are the greatest
reductions retaining respectively every value and every jet. Every
quotient satisfying (QS5.3) has the unique induced map \[
Q/K\longrightarrow Q/J_{\mathrm{jet}},\qquad
q+K\longmapsto q+J_{\mathrm{jet}}.
\tag{QS5.4}
\] Each original \(j_{\rho,m_\rho}\) is still onto from this terminal
observation quotient, since it was onto from \(Q\). Every actual
character and original primary length survives.

For value preservation alone, the nonzero maps \(F_\bullet(\rho)\)
remain onto \(\mathbb C\). On an equivariant quotient,
\(j_{\rho,m_\rho}(K)\) is \(W_p\)-invariant, hence is an ideal by QS2.
Value preservation forces that ideal into \((t_\rho)\). The induced
local quotient therefore has positive length: jets may shorten, but the
character cannot disappear.

No equality \(J_{\mathrm{jet}}=0\), or unsupplied spectral-synthesis
theorem, is asserted. Removing any extra vectors in this intersection
still changes no zero character. IH6 already provides a complete
joint-grid receiver with kernel \(J_{\mathrm{val}}\), giving a concrete
injective realization of \(Q/J_{\mathrm{val}}\).

Preserving the original meromorphic zeta function has an additional
elementary consequence: equality in its original variable fixes all zero
germs and their orders. It cannot simultaneously assert a changed zero
set of that same function. A new conclusion about those zeros' positions
requires an estimate rather than another invocation of uniqueness.

\subsection{QS6. Actual tensor powers with every nilpotent
retained}\label{qs6.-actual-tensor-powers-with-every-nilpotent-retained}

The \(k\)-fold tensor power of one block is exactly \[
A_{\rho,m}^{\otimes k}
=\mathbb C[t_1,\ldots,t_k]/(t_1^m,\ldots,t_k^m).
\tag{QS6.1}
\] With \(N=t_1+\cdots+t_k\), its diagonal prime multiplier is \[
w_{p,\rho,m}^{\otimes k}
=p^{k\rho}\prod_{i=1}^k\left(\sum_{j=0}^{m-1}
\frac{(\log p)^jt_i^j}{j!}\right)
=p^{k\rho}\sum_{j=0}^{k(m-1)}\frac{(\log p)^jN^j}{j!}.
\tag{QS6.2}
\] The multinomial identity proves this equality in the stated quotient.
Its dimension is \(m^k\). The nilpotence order of \(N\) is exactly
\(k(m-1)+1\): larger total degree forces an exponent at least \(m\),
while \[
N^{k(m-1)}
=\frac{(k(m-1))!}{((m-1)!)^k}
\,\prod_{i=1}^k t_i^{m-1}\ne0.
\tag{QS6.3}
\] For \(m=1\), this says \(N^0=1\) and \(N=0\).

For every repetition \(r\ge1\), \[
\operatorname{Tr}((W_p^{\otimes k})^r)=m^kp^{kr\rho},
\qquad
\det(XI-W_p^{\otimes k})=(X-p^{k\rho})^{m^k}.
\tag{QS6.4}
\] Every positive-degree term is strictly triangular in a degree-ordered
monomial basis. This proves both identities without deleting nilpotents
from the actual operator.

For the full four-point algebra \(B_\rho\), distribute the tensor
product: \[
B_\rho^{\otimes k}
=\prod_{(\omega_1,\ldots,\omega_k)\in\mathcal O_\rho^k}
\mathbb C[t_1,\ldots,t_k]/(t_1^m,\ldots,t_k^m).
\tag{QS6.5}
\] On its indicated factor the multiplier is exactly \[
p^{\omega_1+\cdots+\omega_k}
\left(\sum_{j=0}^{k(m-1)}\frac{(\log p)^jN^j}{j!}\right).
\tag{QS6.6}
\] All mixed choices of reflected and conjugate characters remain.
Consequently \[
\operatorname{Tr}((W_p^{\otimes k})^r)
=m^k\left(\sum_{\omega\in\mathcal O_\rho}p^{r\omega}\right)^k,
\tag{QS6.7}
\] \[
\det(XI-W_p^{\otimes k})
=\prod_{(\omega_1,\ldots,\omega_k)\in\mathcal O_\rho^k}
\left(X-p^{\omega_1+\cdots+\omega_k}\right)^{m^k}.
\tag{QS6.8}
\] Coincident eigenvalues add their positive integer multiplicities;
they do not cancel.

\subsection{QS7. The tensor exponent error is computed, not assumed
bounded}\label{qs7.-the-tensor-exponent-error-is-computed-not-assumed-bounded}

For a finite matrix \(A\) with nonzero distinct eigenvalues
\(\lambda_j\) and positive integer algebraic multiplicities \(d_j\), \[
\sum_{r\ge0}\operatorname{Tr}(A^r)z^r
=\sum_j\frac{d_j}{1-\lambda_jz}
\tag{QS7.1}
\] near zero. Distinct denominators have distinct poles and nonzero
residues. The radius of this series is \(1/\max_j|\lambda_j|\), so
Cauchy's coefficient-radius formula gives \[
\limsup_{r\to\infty}|\operatorname{Tr}(A^r)|^{1/r}
=\max_j|\lambda_j|.
\tag{QS7.2}
\] This handles cancellation among trace terms and all multiplicities;
it requires no lower bound for each individual repetition.

Apply this to the actual operator on \(B_\rho^{\otimes k}\). Set \[
b=\max(\beta,1-\beta),\qquad
\delta=b-\tfrac12=|\beta-\tfrac12|.
\] All tuples have real exponent at most \(kb\), and repeating a
character with real part \(b\) attains it. Therefore \[
L_{p,k}
=\limsup_{r\to\infty}
\frac{\log|\operatorname{Tr}((W_p^{\otimes k})^r)|}{r\log p}
=kb,
\] \[
\boxed{E_{p,k}:=L_{p,k}-\frac{k}{2}=k\delta.}
\tag{QS7.3}
\] The logarithm at a zero trace term is \(-\infty\), which does not
affect (QS7.2). Every prime action and every repetition is retained in
this definition.

The actual estimate supplied by the critical strip is \[
|\operatorname{Tr}((W_p^{\otimes k})^r)|
\le(4m)^kp^{rk},\qquad r,k\ge1.
\tag{QS7.4}
\] Every one of the four summands in (QS6.7) has modulus at most
\(p^r\). This proves (QS7.4). Its error relative to \(k/2\) is \(k/2\),
growing with \(k\), so it supplies no amplification to zero.

Among the calculated nonnegative sequences \(E_{p,k}=k\delta\), the ones
with bounded supremum are exactly the zero sequence. More generally the
sequences with \(E_{p,k}/k\to0\) are exactly the zero sequence. These
follow directly from \(\delta\ge0\). They classify the error sequences;
they do not assert that the actual receiver has either bound.

For comparison with the numerical form of a Deligne amplification, an
estimate shaped as \[
|\operatorname{Tr}((W_p^{\otimes k})^r)|
\le C_k r^{d_k}p^{r(k/2+c)}
\tag{QS7.5}
\] has exponent at most \(k/2+c\), because for fixed \(k\) its \(r\)-th
root loses the factors \(C_k^{1/r}r^{d_k/r}\). Its exponent error is the
fixed \(c\), independent of \(k\). Comparing the exact (QS7.3) gives
\(k\delta\le c\), and the preceding sequence classification removes
\(\delta>0\).

Equation (QS7.5) is a comparison of estimate shapes. It is not proved
for the programme in this note and is not supplied as an assumption
completing RH. The proved results are the actual exponent (QS7.3),
actual bound (QS7.4), and full quotient classification. Every
positive-length reflection-compatible quotient retains the same
\(b,\delta\); shortening multiplicities changes finite trace
coefficients but not the exponential rate.

The tensor product here is the complex-linear receiving tensor product
proved in QS6. No tensor comparison from a source cohomology is assumed
by changing notation, and no operation on source \(\tau\) appears in its
exponent computation.

\subsection{QS8. Every support label is retained with the exact
kernel}\label{qs8.-every-support-label-is-retained-with-the-exact-kernel}

Let \(C_{\mathrm{coeff}}\) be the programme's formal receiving
coefficient algebra, with separate recorded symbols \(U=[\tau]\),
\(f=[1]\), and let \(W_L\) be the vector space on all specified support
labels. Put \(V=C_{\mathrm{coeff}}\otimes_\mathbb C W_L\). Their vector
operations define no addition on source \(\tau\).

Use exactly \[
q\otimes\operatorname{id}_V,\qquad
j_{\rho,m}\otimes\operatorname{id}_V.
\tag{QS8.1}
\] Choose a basis of \(V\). Every algebraic tensor is a finite sum
\(\sum_\alpha x_\alpha\otimes v_\alpha\) in independent basis elements.
Its image is zero precisely when each \(q(x_\alpha)\) is zero. Hence \[
\ker(q\otimes\operatorname{id}_V)=(\ker q)\otimes V.
\tag{QS8.2}
\] Every named support coordinate remains the same coordinate. A
surviving character survives in every nonzero coefficient coordinate; no
supported label is identified with source absence by these maps.

This is the exact identity-on-labels extension. It does not assert that
every possible quotient of a larger mixed coefficient category preserves
labels. Such a quotient needs its own definition.

\subsection{QS9. Consequence for the proposed global quotient
route}\label{qs9.-consequence-for-the-proposed-global-quotient-route}

The canonical original-zeta reconstruction fixes the arithmetic measure
and original meromorphic function. The quotient classification now
gives:

\begin{enumerate}
\def\labelenumi{\arabic{enumi}.}
\tightlist
\item
  Killing a block removes its zero observation; it is detectable.
\item
  Shortening a block retains the character and loses specified higher
  derivatives; it does not retain complete original jets.
\item
  Removing the global joint observation kernel retains every
  observation, including any off-critical one.
\item
  Equivariant isomorphism for fixed prime actions preserves every
  character; \(2,3\) already separate them.
\item
  Tensor powers of a surviving off-critical orbit have exact exponent
  error \(k|\beta-\tfrac12|\), with all mixed sectors and nilpotents
  present.
\end{enumerate}

A quotient invisible to all retained observations therefore cannot make
an observed off-critical character critical. It can remove unobserved
receiver data. This proves the distinction through the actual quotient
maps and kernels.

No proof or disproof of RH is asserted. The result identifies the
invariant that a geometric weight estimate must control and calculates
why invisible quotienting or meromorphic uniqueness alone cannot change
it.

\subsection{Reading and source
identity}\label{reading-and-source-identity}

Read completely for this derivation:

\begin{itemize}
\tightlist
\item
  IH1--IH8, original local programme source:
  \nolinkurl{INTEGRAL_HISTORY_AND_PURITY.md}
\item
  PL1--PL8, original local programme source:
  \nolinkurl{PRIME_CLOCK_LOOPS_AND_Z1.md}
\end{itemize}

QS2--QS8 supply complete finite-algebra, quotient, and tensor proofs. No
historical novelty is asserted for these elementary mechanisms. This
independent note has not read Deligne's full source; the coordinating
task is undertaking the human-source reconstruction and must supply
precise theorem locators. No alleged Deligne theorem is used as an input
here.
